\begin{lemma}
Let $\Omega\subseteq\mathbb{C}$ be open, let $j\in\{0,1\}$, write $e_0=1$ and
$e_1=i$, and let $u,D\colon\mathbb{C}\to\mathbb{C}$ satisfy
\begin{enumerate}
\item[(H1)] $u\in L^{1}_{\mathrm{loc}}(\Omega)$ and
  $D\in L^{1}_{\mathrm{loc}}(\Omega)$, in the sense of \noderef{MemLpLocOn};
\item[(H2)] $D$ is the distributional $j$-th partial derivative of $u$ on
  $\Omega$, in the sense of \noderef{IsWeakPartialOn}.
\end{enumerate}
Let $K\subseteq\Omega$ be compact.  Then there is a sequence $(t_l)_{l\ge1}$ of
nonzero real numbers with $t_l\to0$ such that
\[
  \frac{u(z+t_le_j)-u(z)}{t_l}\longrightarrow D(z)\qquad(l\to\infty)
\]
for almost every $z\in K$.

This is the statement that the distributional derivative is a genuine limit of
difference quotients of $u$, almost everywhere, and it is what makes the
distributional derivative comparable with a classical one.  The paper uses it
implicitly throughout Sections 4 to 6, where the $W^{1,2}$ derivatives of $H$
and of the inverse $\chi$ are manipulated as classical derivatives.

Two features of the statement are deliberate and are exactly what the
argument supplies.  First, the increments form a sequence rather than a
continuous parameter: the difference quotients converge to $D$ in
$L^{1}(K)$ as $t\to0$, and almost-everywhere convergence is obtained from
that only along a subsequence.  A full statement, with $t\to0$ continuously
and one exceptional null set, is false in general, because the exceptional
set depends on the increment.  Second, the conclusion is confined to a
compact $K$, since the translate $z+t e_j$ must remain in $\Omega$.  Neither
restriction costs anything where the statement is used: at a point where $u$
is differentiable, every null sequence of increments gives the classical
derivative, and $\Omega$ is a countable union of compact sets.
\end{lemma}
\begin{proof}
Write $d\sigma$ for plane Lebesgue measure.

\emph{Step 1: a compact collar and the truncation of $D$.}  Since $K$ is
compact, $\Omega$ is open and $K\subseteq\Omega$, there is $\delta_1>0$ with
\[
  K^{\delta_1}:=\{z:\operatorname{dist}(z,K)\le\delta_1\}\subseteq\Omega ;
\]
either $\Omega=\mathbb{C}$, where any $\delta_1>0$ serves, or
$\operatorname{dist}(K,\mathbb{C}-\Omega)>0$ because $K$ is compact and
$\mathbb{C}-\Omega$ is closed, nonempty and disjoint from $K$, and then any
$\delta_1$ below that distance serves.  The set $K^{\delta_1}$ is closed and
bounded, hence compact.  Put $V=\mathbf{1}_{K^{\delta_1}}D$.  By (H1) applied
to the compact set $K^{\delta_1}\subseteq\Omega$ and \noderef{MemLpLocOn},
$V\in L^{1}(\mathbb{C})$; by the same hypothesis applied to $K^{\delta_1}$,
$u$ is almost everywhere strongly measurable on $K^{\delta_1}$, so for each
fixed $t$ the difference quotient below is almost everywhere strongly
measurable on $K$.

\emph{Step 2: the line-integral identity.}  By \noderef{WeakPartialLineIntegral}
applied to $\Omega$, $j$, $u$, $D$ and $K$ there is $\delta_2>0$ such that for
every real $t$ with $|t|\le\delta_2$,
\begin{equation}\tag{1}
  u(z+te_j)-u(z)=\int_0^tD(z+se_j)\,ds
  \qquad\text{for almost every }z\in K .
\end{equation}
Put $\delta=\min\{\delta_1,\delta_2\}>0$.

\emph{Step 3: the difference quotients converge in $L^{1}(K)$.}  For
$0<|t|\le\delta$ put
\[
  \Delta_t(z)=\frac{u(z+te_j)-u(z)}{t}\qquad(z\in\mathbb{C}).
\]
Since $\frac1t\int_0^tD(z)\,ds=D(z)$ for every $z$, (1) gives, for almost
every $z\in K$,
\begin{equation}\tag{2}
  \Delta_t(z)-D(z)=\frac1t\int_0^t\bigl[D(z+se_j)-D(z)\bigr]\,ds .
\end{equation}
Let $z\in K$ and $|s|\le\delta$.  Then
$\operatorname{dist}(z+se_j,K)\le|s|\le\delta_1$, so $z+se_j\in K^{\delta_1}$
and hence $D(z+se_j)=V(z+se_j)$; and $z\in K\subseteq K^{\delta_1}$, so
$D(z)=V(z)$.  Therefore the bracket in (2) may be replaced by
$V(z+se_j)-V(z)$.  Writing $I_t$ for the interval between $0$ and $t$, whose
length is $|t|$, and noting that $(z,s)\mapsto V(z+se_j)-V(z)$ is measurable
on $\mathbb{C}\times I_t$, being built from the measurable $V$ and the
continuous map $(z,s)\mapsto z+se_j$, we get from (2), the triangle
inequality for integrals, Tonelli's theorem and $K\subseteq\mathbb{C}$ that
\[
  \int_K\bigl|\Delta_t-D\bigr|\,d\sigma
  \le\frac1{|t|}\int_{I_t}\Bigl(\int_{\mathbb{C}}
      \bigl|V(z+se_j)-V(z)\bigr|\,d\sigma(z)\Bigr)ds
  \le\sup_{|s|\le|t|}\bigl\|V(\cdot+se_j)-V\bigr\|_{L^{1}(\mathbb{C})} ,
\]
the last step because the inner integral is at most that supremum for every
$s\in I_t$ and $|I_t|=|t|$.  Translation is continuous in
$L^{1}(\mathbb{C})$, so $\|V(\cdot+se_j)-V\|_{L^{1}(\mathbb{C})}\to0$ as
$s\to0$, and therefore the right-hand side tends to $0$ as $t\to0$.  Hence
\begin{equation}\tag{3}
  \int_K\bigl|\Delta_t-D\bigr|\,d\sigma\longrightarrow0
  \qquad\text{as }t\to0 ,
\end{equation}
that is, $\Delta_t\to D$ in $L^{1}(K)$.

\emph{Step 4: extraction of an almost everywhere convergent sequence.}  Put
$\tau_l=\delta/l$ for $l\ge1$.  Then $\tau_l\neq0$, $|\tau_l|\le\delta$ and
$\tau_l\to0$, so by (3) we have $\int_K|\Delta_{\tau_l}-D|\,d\sigma\to0$,
i.e.\ $\Delta_{\tau_l}\to D$ in $L^{1}(K)$.  Convergence in $L^{1}(K)$
implies convergence in measure on $K$: for $\eta>0$,
\[
  \sigma\bigl(\{z\in K:|\Delta_{\tau_l}(z)-D(z)|\ge\eta\}\bigr)
  \le\frac1\eta\int_K\bigl|\Delta_{\tau_l}-D\bigr|\,d\sigma
  \longrightarrow0 .
\]
A sequence of almost everywhere strongly measurable functions converging in
measure has a subsequence converging almost everywhere; let
$l_1<l_2<\cdots$ index such a subsequence and put $t_i=\tau_{l_i}$.  Then
each $t_i$ is nonzero, $t_i\to0$ because $l_i\to\infty$, and
$\Delta_{t_i}\to D$ almost everywhere on $K$, which is precisely the
assertion.
\end{proof}
